\section{Phase fringes, arithmetic means, and sampling}
\label{sec:readout}

The source already contains a direct encoding of the mean of a bounded
observable in Eqs.~(32)--(34). A separate phase-reference procedure in
Appendix~D.3 requires a different analysis. We take its probability in
source Eq.~(D6) as the starting point, avoiding the missing affine
probability conversion in source Eq.~(37).

\subsection{The exact statistic measured by a reference scan}
Let $\Theta$ be the real observable encoded in the phase and let $a$ be
the reference angle. On either a finite probability grid or a continuum
probability space, define
\begin{equation}
 P(a):=\E{\cos^2(\Theta-a)},\qquad z:=\E{e^{2i\Theta}}.
 \label{eq:phase-statistic}
\end{equation}
The finite-grid formula is exact, with the probabilities summing to one;
no Riemann-sum approximation is required for the following statements.

\begin{proposition}[Exact phase-reference fringe]
\label{prop:fringe}
For every real $a$,
\begin{equation}
 P(a)=\frac12+\frac12\operatorname{Re}(e^{-2ia}z).
 \label{eq:exact-fringe}
\end{equation}
If $z\ne0$, its maximizing reference satisfies
$a_*=\tfrac12\arg z\pmod\pi$, and its contrast is $|z|$.
This circular mean generally differs from $\E{\Theta}$.
\end{proposition}

\begin{proof}
Use $\cos^2u=(1+\cos2u)/2$ and linearity of expectation.
Writing $z=|z|e^{i\varphi}$ gives
$P(a)=1/2+(|z|/2)\cos(\varphi-2a)$, proving the assertion.
\end{proof}

Source Eq.~(D8) correctly displays variance as a leading reduction of the
fringe amplitude. The paragraph after source Eq.~(D10), p.~15, then
introduces $\cos(2\theta_{\rm err})=1-2\operatorname{Var}(\Theta)$ and
uses an additive shift in a squared cosine. Multiplying a cosine's
amplitude and shifting its argument are different operations: a fixed
shift generally introduces an odd component about the original peak.
Equation~\eqref{eq:exact-fringe} identifies the relevant quantities without
this substitution.

\begin{corollary}[Positive variance without phase bias]
\label{cor:symmetric-no-bias}
If $Y:=\Theta-\mu$ has a symmetric distribution and
$C:=\E{\cos2Y}>0$, then
\begin{equation}
 P(a)=\frac12+\frac C2\cos[2(\mu-a)],\qquad a_*=\mu\pmod\pi.
 \label{eq:symmetric-fringe}
\end{equation}
The conclusion holds even when $\operatorname{Var}(\Theta)>0$.
\end{corollary}

\begin{proof}
Symmetry gives $\E{\sin2Y}=0$, hence $z=e^{2i\mu}C$.
Apply Proposition~\ref{prop:fringe}.
\end{proof}

For the ideal Gaussian $\Theta\sim N(\mu,\sigma^2)$ on $\R$,
$C=e^{-2\sigma^2}$. A symmetric truncated Gaussian on
$[\mu-b,\mu+b]$ with $b<\pi/4$ also has $C>0$, since $\cos2Y>0$
throughout its support. The same result holds exactly for symmetric
finite-grid weights. Thus the counterexample does not depend on an
unbounded phase-space domain. A reference interval of length less than
$\pi$ containing $\mu$ fixes the relevant branch of the maximum.
Nonzero variance reduces contrast but does not impose a nonzero phase
bias in this family. With increasing samples the maximum can be located
arbitrarily accurately; reduced contrast affects the resources required.

The origin of a phase displacement can also be isolated locally. If
$\Theta=\mu+sZ$, $\E{Z}=0$, and $\E{|Z|^5}<\infty$, Taylor expansion of
$\E{e^{2isZ}}$ on the branch near $1$ yields
\begin{equation}
 a_* =\mu-\frac23s^3\E{Z^3}+O(s^5),\qquad s\to0.
 \label{eq:phase-asymmetry}
\end{equation}
Here $\mu$ is fixed, and $Z$ has a fixed distribution with the stated
moments. The second moment changes the real part at order $s^2$;
the leading imaginary part is $-4s^3\E{Z^3}/3$.
This is a moment expansion of the exact fringe, not a general claim that
all nonlinear inference biases begin at the same order.

\subsection{A fixed-frequency fringe need not identify the arithmetic mean}
The positive-variance counterexample disproves a mandatory phase-bias
floor. Conversely, it does not make the circular mean a valid arithmetic
mean estimator for arbitrary packets. A complete scan of $a$ in
\eqref{eq:exact-fringe} determines only the one complex moment $z$.

\begin{proposition}[Same entire fringe, different arithmetic means]
\label{prop:nonidentifiable}
There exist strictly positive probability weights on four binary-grid
points in $[0,1]$ with identical $P(a)$ for every $a$ and different
arithmetic means.
\end{proposition}

\begin{proof}
Choose $\mu=1/2$, $a=1/8$, $b=1/16$, and the ordered support
$x=(\mu-a,\mu-b,\mu+b,\mu+a)$.
Set $c=\sin(2b)/\sin(2a)$, $v=(-c,1,-1,c)$, and
\begin{equation}
 p_j^{\pm}=\frac14\pm\eta v_j,\qquad 0<\eta<\frac14.
 \label{eq:four-point-weights}
\end{equation}
Since $0<c<1$, these weights are positive. Moreover,
\begin{equation}
 \sum_jv_j=0,\qquad
 \sum_jv_je^{2ix_j}=2ie^{2i\mu}
       \bigl(c\sin(2a)-\sin(2b)\bigr)=0.
 \label{eq:four-point-null}
\end{equation}
The two distributions therefore have the same normalization and the
same $z$. Their means differ by
\begin{equation}
 \sum_j(p_j^+-p_j^-)x_j=4\eta(ac-b)>0.
 \label{eq:four-point-means}
\end{equation}
The strict inequality follows because $\sin t/t$ is decreasing for
$0<t<\pi/2$, so $c>b/a$.
\end{proof}

At $\eta=0.1$ the mean separation is approximately
$1.9659219701\times10^{-4}$. All four support points lie on a
$16$-point binary grid. This is an identifiability statement for the
fixed-frequency reference-fringe family. It is not an impossibility
theorem for algorithms that encode a different function, obtain other
Fourier moments, or measure additional information from the state.

\subsection{Direct mean encoding and its resource assumptions}
Let a state have grid probabilities $p_x$, and let $0\leq f(x)\leq1$.
The Hadamard circuit with diagonal phase
$\vartheta(x)=2\arccos\sqrt{f(x)}$ gives ancilla-zero amplitudes
$\sqrt{p_x}(1+e^{i\vartheta(x)})/2$. Therefore its probability is exactly
\begin{equation}
 p_f=\sum_xp_x\cos^2\!\left(\frac{\vartheta(x)}2\right)
     =\sum_xp_xf(x)=\E{f(X)}.
 \label{eq:direct-mean-probability}
\end{equation}
This is the source's own Eqs.~(32)--(34), read as a finite-grid identity.
Standard amplitude estimation can estimate this encoded probability
with additive accuracy using $O(\epsilon^{-1})$ coherent queries at fixed
success probability \cite{bhmt,aaronsonrall}; more general observable
estimation is also established \cite{rall}. These are query statements
requiring appropriate coherent state preparation, inverse operations,
reflections, and accurate observable encoding. Their gate costs and
state-preparation success costs remain to be counted. No new amplitude
estimation algorithm or end-to-end quantum speedup is claimed here.

\subsection{Two sampling experiments have different variances}
The measurement-cost prefactor in source Sec.~V, p.~6, is written in
terms of the packet covariance. For an ancilla trial $Y\in\{0,1\}$
with success probability $p_f$, however, independent trials give
\begin{equation}
 \operatorname{Var}\!\left(\frac1M\sum_{m=1}^MY_m\right)
 =\frac{p_f(1-p_f)}{M}.
 \label{eq:ancilla-variance}
\end{equation}
If instead one samples the coordinate $X$ directly in the computational
basis, its sample mean has variance $\operatorname{Var}(X)/M$.
These statements concern different recorded random variables.
For $f(X)=X$ and any distribution on $[0,1]$ symmetric about $1/2$,
$p_f=1/2$ and the ancilla variance per shot is $1/4$, irrespective of
the coordinate variance. A narrow Gaussian packet can have arbitrarily
small coordinate variance without making these ancilla outcomes less
random. Source Appendix~D.1 analyzes direct Gaussian samples, while
Appendix~D.2 separately gives the Bernoulli variance factor. The former
cannot be transferred as the prefactor of the latter circuit.

The inverse-square accuracy exponent for ordinary independent sampling
survives this correction. For example, Hoeffding's inequality for the
ancilla sample mean gives failure probability at most
$2e^{-2M\epsilon^2}$. The correction concerns the measured variable and
its variance prefactor, rather than a refutation of standard sampling or
amplitude-estimation scaling. Finally, estimating the actual packet mean
accurately does not remove the separate difference between that mean
and the nonlinear trajectory launched at its initial center.
